%! Graphing Radical Functions \& Transformations — Unit 6, Lesson 6.4 — Guided Notes (Answer Key)
\documentclass[10pt]{article}
\usepackage{algebra2-article}
\usepackage{algebra2-key}   % pulls in -boxes; do NOT also load -boxes
\newcommand{\TallMath}[1]{$\displaystyle #1\rule[-1.4em]{0pt}{3.2em}$}
% In-formula answer slot (\ans is text-mode and must never appear inside $...$).
\newcommand{\mans}[1]{{\color{keyred}\mathbf{#1}}}
% Vocab answer for the key: bold forest term, then the filled definition on a write-line.
% The leading and trailing \par are required: \ansline ends with \dotfill but does NOT end the
% paragraph, so without them each term label gets pulled onto the previous answer's line.
\newcommand{\vocabans}[2]{%
  \par\noindent\textbf{\textcolor{forest}{#1:}}\\[1pt]\ansline{#2}\par}
\newcommand{\lgrid}[4]{%
  \draw[step=1, linegray!40, very thin] (#1,#3) grid (#2,#4);
  \draw[->, semithick, charcoal] (#1,0)--(#2,0) node[below right, font=\scriptsize]{$x$};
  \draw[->, semithick, charcoal] (0,#3)--(0,#4) node[left, font=\scriptsize]{$y$};}
% Cube root a*cbrt(x-h)+k: pgfmath has no cbrt, so plot the two branches separately.
% #1=a  #2=h  #3=k  #4=xmin  #5=xmax
\newcommand{\cbrtplot}[5]{%
  \draw[very thick, forest, domain=#2:#5, samples=70, smooth]
    plot (\x,{(#1)*pow((\x)-(#2),1/3)+(#3)});
  \draw[very thick, forest, domain=#4:#2, samples=70, smooth]
    plot (\x,{-(#1)*pow((#2)-(\x),1/3)+(#3)});}

\begin{document}

\pageheader{Unit 6, Lesson 6.4}{Guided Notes --- Answer Key}
\namedateperiod

\vspace{0.08in}

\begin{objectivebox}
\textbf{By the end of this lesson, I will be able to\dots}
\begin{itemize}\setlength{\itemsep}{2pt}
  \item graph $y=a\sqrt{x-h}+k$ and $y=a\sqrt[3]{x-h}+k$ from the equation by moving the parent's \emph{anchor point} to $(h,k)$ and applying $a$, and read off the domain and range without solving anything;
  \item write the \emph{equation} of a square root or cube root function from its graph, using the anchor point for $h$ and $k$ and one more point for $a$;
  \item explain why, for a radical function, changing the \emph{input} by a factor and changing the \emph{output} by a factor produce the same graph.
\end{itemize}
\end{objectivebox}

\vspace{0.05in}

\begin{vocabbox}
\small
Fill in each term as we build it together.
\par\vspace{2pt}   % \par is required: \vocabans's \noindent is a no-op mid-paragraph
\vocabans{Transformation form of a radical function}{$y=a\sqrt{x-h}+k$ or $y=a\sqrt[3]{x-h}+k$ --- the form in which every transformation from the parent can be read straight off the equation}
\vocabans{Anchor point}{the one point a radical graph is built around, always $(h,k)$: the \emph{endpoint} for a square root, the \emph{inflection point} for a cube root. Find it and the rest of the graph follows}
\vocabans{Vertical stretch / vertical compression ($a$)}{multiplying every output by $a$. $|a|>1$ stretches the graph away from the $x$-axis (steeper); $0<|a|<1$ compresses it toward the $x$-axis (flatter). The anchor never moves}
\vocabans{Reflection across the $x$-axis}{what a \emph{negative} $a$ does --- every output changes sign, so the arm leaves the anchor going down and the anchor becomes an absolute \emph{maximum} instead of a minimum}
\vocabans{Inflection point (the cube root's anchor)}{the point $(h,k)$ where a cube root graph flattens out and switches the direction it bends. It is \emph{not} an extremum --- the graph keeps going in both directions}
\vocabans{Horizontal compression, $f(kx)$}{multiplying every \emph{input} by $k$, sliding each point toward the $y$-axis. For a radical it cannot be told apart from a vertical stretch, because $\sqrt[n]{kx}=\sqrt[n]{k}\cdot\sqrt[n]{x}$}
\end{vocabbox}

\begin{hookbox}
Three graphs. Three equations. They do \emph{not} appear in the same order.
\begin{center}
\begin{tikzpicture}[scale=0.30]
  % I: y = sqrt(x)
  \begin{scope}
    \lgrid{-2}{10}{-4}{4}
    \begin{scope}
      \clip (-2,-4) rectangle (10,4);
      \draw[very thick, forest, domain=0:10, samples=90, smooth] plot (\x,{sqrt(\x)});
    \end{scope}
    \fill[forest] (0,0) circle (6pt);
    \node[charcoal, font=\small] at (4,-5.6) {\textbf{I}};
  \end{scope}
  % II: y = sqrt(x-3)
  \begin{scope}[xshift=15cm]
    \lgrid{-2}{10}{-4}{4}
    \begin{scope}
      \clip (-2,-4) rectangle (10,4);
      \draw[very thick, forest, domain=3:10, samples=90, smooth] plot (\x,{sqrt((\x)-3)});
    \end{scope}
    \fill[forest] (3,0) circle (6pt);
    \node[charcoal, font=\small] at (4,-5.6) {\textbf{II}};
  \end{scope}
  % III: y = sqrt(x) - 3
  \begin{scope}[xshift=30cm]
    \lgrid{-2}{10}{-4}{4}
    \begin{scope}
      \clip (-2,-4) rectangle (10,4);
      \draw[very thick, forest, domain=0:10, samples=90, smooth] plot (\x,{sqrt(\x)-3});
    \end{scope}
    \fill[forest] (0,-3) circle (6pt);
    \node[charcoal, font=\small] at (4,-5.6) {\textbf{III}};
  \end{scope}
\end{tikzpicture}
\end{center}
\vspace{2pt}
Match each equation to its graph: \quad $y=\sqrt{x}-3$ \; is \ans{III} \quad
$y=\sqrt{x-3}$ \; is \ans{II} \quad $y=\sqrt{x}$ \; is \ans{I}
\begin{itemize}\setlength{\itemsep}{1pt}
  \item Two of these equations contain a $3$ and both graphs moved --- but in different directions. What decides? \ansline{Whether the $3$ is \emph{inside} the radical or \emph{outside} it. Inside, it moves the graph horizontally and in the opposite direction to its sign; outside, it moves the graph vertically and in the same direction.}
  \item In Lesson~6.0 you found the domain of $\sqrt{x-3}$ by \emph{solving} $x-3\ge0$. Look at graph~II. Where could you have simply \emph{read} that answer? \ansline{At the endpoint. Graph II starts at $(3,0)$, so the domain is $[3,\infty)$ --- the inequality's answer is just the first coordinate of the point where the graph begins.}
\end{itemize}
Today the domain stops being a calculation. \textbf{Find one point, and the whole graph follows.}
\end{hookbox}

\begin{notesbox}{1. The Anchor Point: $y=\sqrt{x-h}+k$}
Every radical graph is built around a single point. For a square root that point is the
\textbf{endpoint} --- where the graph begins. We call it the \textbf{anchor point}, and here is the
entire rule:
\begin{center}
\begin{tcolorbox}[colback=goldbg, colframe=goldacc, boxrule=0.8pt, arc=1.5mm, width=0.92\linewidth,
    left=3mm, right=3mm, top=2mm, bottom=2mm]
\centering
$y=\sqrt{x-h}+k$ \quad has its anchor point at \quad $(\ \mans{h}\,,\ \mans{k}\ )$,
\quad and the graph is the parent $y=\sqrt{x}$ moved there.
\end{tcolorbox}
\end{center}
The signs behave exactly as they did for the parabola and the absolute value in Warm-Up item~1:
\begin{itemize}[leftmargin=1.5em, itemsep=3pt]
  \item $k$ is \emph{outside} the radical. It moves the graph \ans{vertically}, in the
        \ans{same} direction as its sign.
  \item $h$ is \emph{inside}, with the $x$. It moves the graph \ans{horizontally}, in the
        \ans{opposite} direction. So $\sqrt{x-3}$ moves \ans{right} $3$, and $\sqrt{x+5}$
        moves \ans{left} $5$.
\end{itemize}

\vspace{3pt}
\textbf{Worked example: $f(x)=\sqrt{x-3}+1$.}
\begin{center}
\begin{minipage}[c]{0.52\linewidth}
\centering
\begin{tikzpicture}[scale=0.46]
  \lgrid{-1}{13}{-2}{6}
  \begin{scope}
    \clip (-1,-2) rectangle (13,6);
    \draw[very thick, forest, domain=3:13, samples=100, smooth] plot (\x,{sqrt((\x)-3)+1});
  \end{scope}
  \fill[forest] (3,1) circle (4pt) node[left, font=\scriptsize]{$(3,1)$};
  \fill[charcoal] (4,2) circle (2.6pt);
  \fill[charcoal] (7,3) circle (2.6pt) node[above left, font=\scriptsize]{$(7,3)$};
  \fill[charcoal] (12,4) circle (2.6pt);
\end{tikzpicture}
\end{minipage}
\begin{minipage}[c]{0.44\linewidth}
\small
$h=\mans{3}$, \quad $k=\mans{1}$ \\[4pt]
Anchor point: \ans{$(3,1)$} \\[4pt]
Domain: \ans{$[3,\infty)$} \\[4pt]
Range: \ans{$[1,\infty)$} \\[4pt]
The graph is the parent moved \\ right \ans{3} and up \ans{1}.
\end{minipage}
\end{center}

\vspace{2pt}
\textbf{Why this replaces the inequality.} Solving $x-3\ge0$ gives $x\ge 3$. Reading the anchor
point gives $x$ starting at $3$. Those are not two facts --- they are \ans{the same fact}. Sliding the
graph right by $h$ \emph{is} the act of throwing away every input below $h$.
\begin{center}
\begin{tcolorbox}[colback=forestbg, colframe=forest, boxrule=0.7pt, arc=1.5mm, width=0.92\linewidth,
    left=3mm, right=3mm, top=1.5mm, bottom=1.5mm]
\centering\small
For $y=\sqrt{x-h}+k$: \quad Domain $=[\,h,\infty)$ \quad and \quad Range $=[\,k,\infty)$.
\emph{Both start at the anchor.}
\end{tcolorbox}
\end{center}
\end{notesbox}

\begin{notesbox}{2. The Coefficient $a$: Steeper, Flatter, or Upside Down}
Now put a number in front: $y=a\sqrt{x-h}+k$. The coefficient $a$ multiplies every \emph{output},
so it changes the graph vertically --- and it never moves the anchor point, because at the anchor the
radical equals \ans{0}, and $a\cdot 0 = \mans{0}$.
\begin{center}
\begin{minipage}[c]{0.46\linewidth}
\centering
\begin{tikzpicture}[scale=0.46]
  \lgrid{-1}{10}{-1}{7}
  \begin{scope}
    \clip (-1,-1) rectangle (10,7);
    \draw[very thick, navy, domain=0:10, samples=100, smooth] plot (\x,{2*sqrt(\x)});
    \draw[very thick, forest, domain=0:10, samples=100, smooth] plot (\x,{sqrt(\x)});
    \draw[very thick, goldacc, domain=0:10, samples=100, smooth] plot (\x,{0.5*sqrt(\x)});
  \end{scope}
  \fill[charcoal] (0,0) circle (4pt);
  \node[navy, font=\scriptsize, right] at (6.2,5.4) {$y=2\sqrt{x}$};
  \node[forest, font=\scriptsize, right] at (7.0,2.9) {$y=\sqrt{x}$};
  \node[goldacc, font=\scriptsize, right] at (6.8,1.1) {$y=\tfrac12\sqrt{x}$};
\end{tikzpicture}
\end{minipage}
\hfill
\begin{minipage}[c]{0.50\linewidth}
\small
\begin{itemize}[leftmargin=1.3em, itemsep=4pt]
  \item $|a|>1$: every output is multiplied by more than one, so the graph is a \textbf{vertical
        \ans{stretch}} --- it climbs \ans{faster}.
  \item $0<|a|<1$: a \textbf{vertical \ans{compression}} --- it climbs \ans{slower}.
  \item All three graphs share the same \ans{anchor point} at $(0,0)$, and all three have domain
        \ans{$[0,\infty)$}.
\end{itemize}
\end{minipage}
\end{center}

\vspace{4pt}
\textbf{A negative $a$ turns the arm over.} If $a<0$, every output changes sign: the graph is
\textbf{reflected across the $x$-axis} --- or, more usefully, reflected across the horizontal line
through its own anchor. The arm now leaves the anchor going \emph{down}.
\begin{center}
\begin{minipage}[c]{0.46\linewidth}
\centering
\begin{tikzpicture}[scale=0.46]
  \lgrid{-1}{10}{-4}{7}
  \begin{scope}
    \clip (-1,-4) rectangle (10,7);
    \draw[very thick, forest, domain=0:10, samples=100, smooth] plot (\x,{6-3*sqrt(\x)});
  \end{scope}
  \fill[forest] (0,6) circle (4pt) node[right, font=\scriptsize]{$(0,6)$};
  \fill[charcoal] (1,3) circle (2.6pt);
  \fill[charcoal] (4,0) circle (2.6pt) node[above right, font=\scriptsize]{$(4,0)$};
  \fill[charcoal] (9,-3) circle (2.6pt);
\end{tikzpicture}
\end{minipage}
\hfill
\begin{minipage}[c]{0.50\linewidth}
\small
$m(x)=-3\sqrt{x}+6$ \\[4pt]
$a=\mans{-3}$, \; $h=\mans{0}$, \; $k=\mans{6}$ \\[4pt]
Anchor point: \ans{$(0,6)$} \\[4pt]
Domain: \ans{$[0,\infty)$} \quad Range: \ans{$(-\infty,6]$} \\[4pt]
The anchor is an absolute \ans{maximum}, not a minimum. \\[4pt]
$x$-intercept: \ans{$(4,0)$}
\end{minipage}
\end{center}
\emph{Careful:} the range flips too. When $a>0$ the range is $[k,\infty)$; when $a<0$ it is
\ans{$(-\infty,k]$}.
\end{notesbox}

\begin{notesbox}{3. The Cube Root Version: $y=a\sqrt[3]{x-h}+k$}
Everything above transfers, with \emph{one} change: a cube root has no endpoint, because nothing
stops it. The point at $(h,k)$ is still the anchor --- it is where the graph flattens out and changes
the way it bends --- but it is called the \textbf{inflection point}, and it is not an extremum.
\begin{center}
\begin{minipage}[c]{0.54\linewidth}
\centering
\begin{tikzpicture}[scale=0.38]
  \lgrid{-10}{9}{-5}{2}
  \begin{scope}
    \clip (-10,-5) rectangle (9,2);
    \cbrtplot{1}{-1}{-2}{-10}{9}
  \end{scope}
  \fill[forest] (-1,-2) circle (5pt) node[right, font=\scriptsize]{$(-1,-2)$};
  \fill[charcoal] (-9,-4) circle (3pt);
  \fill[charcoal] (0,-1) circle (3pt);
  \fill[charcoal] (7,0) circle (3pt) node[below right, font=\scriptsize]{$(7,0)$};
\end{tikzpicture}
\end{minipage}
\begin{minipage}[c]{0.42\linewidth}
\small
$g(x)=\sqrt[3]{x+1}-2$ \\[4pt]
$h=\mans{-1}$, \; $k=\mans{-2}$ \\[4pt]
Inflection point: \ans{$(-1,-2)$} \\[4pt]
Domain: \ans{$(-\infty,\infty)$} \\[4pt]
Range: \ans{$(-\infty,\infty)$} \\[4pt]
$x$-intercept: \ans{$(7,0)$} \\[4pt]
$y$-intercept: \ans{$(0,-1)$}
\end{minipage}
\end{center}
\vspace{2pt}
The anchor still tells you where the graph \emph{is}, but it no longer tells you the domain --- the
domain is \ans{all reals} no matter what $h$ is, because an \ans{odd} index refuses nothing.

\vspace{5pt}
\textbf{A bonus only the cube root gets.} Reflect $y=\sqrt[3]{x}$ across the $x$-axis and you get
$y=-\sqrt[3]{x}$. Reflect it across the $y$-axis instead --- replacing $x$ with $-x$ --- and you get
$y=\sqrt[3]{-x}$. Look at what happens:
\begin{center}
\begin{minipage}[c]{0.40\linewidth}
\centering
\begin{tikzpicture}[scale=0.34]
  \lgrid{-8}{8}{-3}{3}
  \begin{scope}
    \clip (-8,-3) rectangle (8,3);
    \cbrtplot{-1}{0}{0}{-8}{8}
  \end{scope}
  \fill[charcoal] (0,0) circle (4pt);
\end{tikzpicture}
\end{minipage}
\begin{minipage}[c]{0.56\linewidth}
\small
This is \textbf{one} graph, not two. Check a point: \\[4pt]
$-\sqrt[3]{8}=\mans{-2}$ \quad and \quad $\sqrt[3]{-8}=\mans{-2}$. \\[4pt]
So $-\sqrt[3]{x}$ and $\sqrt[3]{-x}$ are \ans{equal} for every $x$. \\[4pt]
This is the \textbf{odd symmetry} from Lesson~6.0 doing work: flipping a cube root vertically and
flipping it horizontally are the \emph{same} transformation.
\end{minipage}
\end{center}
\emph{Contrast:} $-\sqrt{x}$ (opens right, going down) and $\sqrt{-x}$ (opens left, going up) are
\ans{completely different} graphs. The square root has \ans{no} symmetry to trade on.
\end{notesbox}

\begin{notesbox}{4. Backwards: From a Graph to an Equation}
This is the harder direction, and it is on the SOL. Only three numbers are unknown, and the graph
hands them to you in a fixed order.
\begin{center}
\begin{minipage}[c]{0.48\linewidth}
\centering
\begin{tikzpicture}[scale=0.44]
  \lgrid{-4}{9}{-2}{9}
  \begin{scope}
    \clip (-4,-2) rectangle (9,9);
    \draw[very thick, forest, domain=-2:9, samples=100, smooth] plot (\x,{3*sqrt((\x)+2)-1});
  \end{scope}
  \fill[forest] (-2,-1) circle (4pt) node[below right, font=\scriptsize]{$(-2,-1)$};
  \fill[charcoal] (2,5) circle (2.6pt) node[right, font=\scriptsize]{$(2,5)$};
  \fill[charcoal] (7,8) circle (2.6pt) node[right, font=\scriptsize]{$(7,8)$};
\end{tikzpicture}
\end{minipage}
\begin{minipage}[c]{0.48\linewidth}
\small
\textbf{Step 1 --- read the anchor.} It is at \ans{$(-2,-1)$}, so \\[2pt]
$h=\mans{-2}$ and $k=\mans{-1}$. Write down what you know:
\[ y=a\sqrt{x-(\mans{-2})}+(\mans{-1}) \]
\textbf{Step 2 --- pick one more lattice point.} Use $(2,5)$. \\[2pt]
\textbf{Step 3 --- substitute and solve for $a$.}
\end{minipage}
\end{center}
\vspace{-4pt}
\begin{center}
\renewcommand{\arraystretch}{1.7}
\begin{tabular}{@{}l l@{}}
$5=a\sqrt{2+2}-1$ & substitute the point \\
$5=a\cdot\mans{2}-1$ & the radical is a \emph{number} --- evaluate it first \\
$\mans{6}=2a$ & add $1$ to both sides \\
$a=\mans{3}$ & \\
\end{tabular}
\end{center}
\textbf{Equation:} $y=$ \ans{$3\sqrt{x+2}-1$} \quad \textbf{Check} with the other point $(7,8)$:
$3\sqrt{7+2}-1=3\cdot\mans{3}-1=\mans{8}$. \; \checkmark
\vspace{3pt}

\emph{The order matters.} If you try to find $a$ first you have two unknowns in one equation. The
anchor is free information --- take it first, every time.
\end{notesbox}

\begin{notesbox}{5. When $h$ Is Hiding --- and the Surprise from Warm-Up Item 3}
Sometimes the radicand is not in the form $x-h$. \textbf{Factor out the coefficient of $x$ first.}
\[
  y=\sqrt{4x-8} \;=\; \sqrt{4(x-\mans{2})} \;=\; \sqrt{4}\cdot\sqrt{x-2}
    \;=\; \mans{2}\sqrt{x-2}
\]
Read the last form: anchor at \ans{$(2,0)$}, and a vertical stretch by \ans{2}.
\begin{center}
\begin{minipage}[c]{0.46\linewidth}
\centering
\begin{tikzpicture}[scale=0.44]
  \lgrid{-1}{12}{-1}{7}
  \begin{scope}
    \clip (-1,-1) rectangle (12,7);
    \draw[very thick, forest, domain=2:12, samples=100, smooth] plot (\x,{2*sqrt((\x)-2)});
  \end{scope}
  \fill[forest] (2,0) circle (4pt) node[below right, font=\scriptsize]{$(2,0)$};
  \fill[charcoal] (3,2) circle (2.6pt);
  \fill[charcoal] (6,4) circle (2.6pt) node[above left, font=\scriptsize]{$(6,4)$};
  \fill[charcoal] (11,6) circle (2.6pt);
\end{tikzpicture}
\end{minipage}
\begin{minipage}[c]{0.50\linewidth}
\small
Domain: \ans{$[2,\infty)$} \quad Range: \ans{$[0,\infty)$} \\[5pt]
\emph{Never} read $h$ off $\sqrt{4x-8}$ directly. The graph does \textbf{not} start at $x=8$.
Factor, and the truth appears: it starts at $x=\mans{2}$. \\[5pt]
(Or check it the old way: $4x-8\ge0$ gives $x\ge\mans{2}$. The two methods must agree ---
they always do.)
\end{minipage}
\end{center}

\vspace{4pt}
\textbf{Now the strange part.} The SOL lists four transformations: $f(x)+k$, $f(x+k)$, $kf(x)$, and
$f(kx)$. That last one, $f(kx)$, multiplies the \emph{input} --- it squeezes the graph
\emph{horizontally}. Watch what it does to a square root:
\begin{center}
\renewcommand{\arraystretch}{1.5}
\begin{tabularx}{0.94\linewidth}{@{}X X@{}}
\hline
\textbf{Squeeze the input: $f(9x)$} & \textbf{Stretch the output: $3f(x)$} \\
\hline
$y=\sqrt{9x}$ --- a \emph{horizontal} compression: each point slides in toward the $y$-axis, to
one \ans{ninth} of its old $x$-value. &
$y=$ \ans{$3\sqrt{x}$} --- a \emph{vertical} stretch: each output is multiplied by
\ans{3}. \\
\hline
\end{tabularx}
\end{center}
These are the \emph{same expression} (Warm-Up item~3), so they are the same graph. Fill in the table
to see it:
\begin{center}
\renewcommand{\arraystretch}{1.3}
\begin{tabular}{|c|c|c|c|c|}
\hline
$x$ & $0$ & $1$ & $4$ & $9$ \\ \hline
$\sqrt{9x}$ & \ans{$0$} & \ans{$3$} & \ans{$6$} & \ans{$9$} \\ \hline
$3\sqrt{x}$ & \ans{$0$} & \ans{$3$} & \ans{$6$} & \ans{$9$} \\ \hline
\end{tabular}
\end{center}
\textbf{Why does this happen?} Because the product property lets a constant walk straight out from
under the radical: $\sqrt{kx}=\sqrt{k}\cdot\sqrt{x}$. So for a radical function --- and only because
of that property --- a horizontal size change \emph{is} a vertical size change wearing a different
coat. There is really only \ans{one} size dial, and you may describe it either way.
\begin{center}
\begin{tcolorbox}[colback=forestbg, colframe=forest, boxrule=0.7pt, arc=1.5mm, width=0.94\linewidth,
    left=3mm, right=3mm, top=1.5mm, bottom=1.5mm]
\small For a cube root the same thing happens with the cube root of $k$:
$\sqrt[3]{8x}=\mans{2}\sqrt[3]{x}$. \\
\emph{This is not true of every family.} Try it on $y=x+1$: compressing the input gives
$\,9x+1$, but stretching the output gives $\,9x+\mans{9}$. Not the same. The
\ans{$+1$} is what breaks it.
\end{tcolorbox}
\end{center}
\end{notesbox}

\begin{practicebox}
\textbf{1.} The graph below is $p(x)=-2\sqrt{x+1}+4$.
\begin{center}
\begin{tikzpicture}[scale=0.46]
  \lgrid{-3}{11}{-4}{6}
  \begin{scope}
    \clip (-3,-4) rectangle (11,6);
    \draw[very thick, forest, domain=-1:11, samples=100, smooth] plot (\x,{4-2*sqrt((\x)+1)});
  \end{scope}
  \fill[forest] (-1,4) circle (4pt) node[above right, font=\scriptsize]{$(-1,4)$};
  \fill[charcoal] (0,2) circle (2.6pt);
  \fill[charcoal] (3,0) circle (2.6pt) node[above right, font=\scriptsize]{$(3,0)$};
  \fill[charcoal] (8,-2) circle (2.6pt);
\end{tikzpicture}
\end{center}
\begin{enumerate}[leftmargin=1.7em, itemsep=3pt]
  \item $a=\mans{-2}$, \; $h=\mans{-1}$, \; $k=\mans{4}$ \quad Anchor point:
        \ans{$(-1,4)$}
  \item \textbf{Domain:} \ans{$[-1,\infty)$} \quad \textbf{Range:} \ans{$(-\infty,4]$}
  \item \textbf{$x$-intercept:} \ans{$(3,0)$} \quad \textbf{$y$-intercept:} \ans{$(0,2)$}
  \item Describe the three transformations from $y=\sqrt{x}$, in words. \ansline{Left $1$ (from
        $x+1$); a vertical stretch by a factor of $2$ together with a reflection across the
        $x$-axis (from $a=-2$); then up $4$.}
  \item The anchor is an absolute \ans{maximum}. Which part of the equation told you that, before
        you looked at the graph? \ans{the negative sign on $a=-2$}
\end{enumerate}

\vspace{4pt}
\textbf{2.} Write the equation of the cube root function graphed below, in the form
$y=a\sqrt[3]{x-h}+k$.
\begin{center}
\begin{tikzpicture}[scale=0.40]
  \lgrid{-8}{10}{-6}{4}
  \begin{scope}
    \clip (-8,-6) rectangle (10,4);
    \cbrtplot{2}{1}{-1}{-8}{10}
  \end{scope}
  \fill[forest] (1,-1) circle (5pt) node[right, font=\scriptsize]{$(1,-1)$};
  \fill[charcoal] (9,3) circle (3pt) node[below right, font=\scriptsize]{$(9,3)$};
  \fill[charcoal] (-7,-5) circle (3pt);
\end{tikzpicture}
\end{center}
Inflection point \ans{$(1,-1)$} $\Rightarrow$ $h=\mans{1}$, $k=\mans{-1}$. \quad
Substitute $(9,3)$: \; $3=a\sqrt[3]{\mans{8}}-1$, so $a=\mans{2}$. \\[3pt]
\textbf{Equation:} $y=$ \ans{$2\sqrt[3]{x-1}-1$}
\end{practicebox}

\begin{teachernote}
\small
\textbf{Pacing.} Sections 1--3 are one continuous idea and should move fast --- students have been
doing $h$ and $k$ since Unit~2, and the only genuinely new word is \emph{anchor point}. Budget the
saved time for Sections 4 and 5, which are where the lesson is won or lost.

\textbf{Section 1 is the argument of the day.} Do not let ``domain $=[h,\infty)$'' land as a new rule
to memorize. It is the \emph{old} rule, seen from the other side: the inequality
$x-h\ge0$ and the horizontal shift by $h$ are one event. If a student is shaky, have them do both and
compare --- they will never disagree.

\textbf{Section 3's reflection pair} is the compare-and-contrast item the SOL asks for
(\textbf{A2.F.1e}), and it lands only if you draw the square-root contrast beside it. $-\sqrt x$ and
$\sqrt{-x}$ occupy different quadrants; $-\sqrt[3]{x}$ and $\sqrt[3]{-x}$ are the same curve. The
reason is Lesson 6.0's odd symmetry, so name it.

\textbf{Section 4 must be practiced in order.} The single most common failure is students trying to
solve for $a$ and $h$ at once. Say it as a rule: \emph{the anchor is free; take it first.} Also
insist on the check with the second point --- it costs ten seconds and catches every sign error.

\textbf{Section 5 is the hard one, and it has two halves.} The first (factor before reading $h$) is
mechanical and will be on every assessment: $\sqrt{4x-8}$ does \emph{not} start at $8$. The second
(the $f(kx)$ collapse) is conceptual and is the answer to the vote you took in the Warm-Up. Do not
let it become ``horizontal and vertical stretches are always the same'' --- the $y=x+1$
counterexample in the green box exists to stop exactly that, and it should be worked aloud.
\end{teachernote}

\end{document}
